%% SIGGRAPH (ACM Transactions on Graphics) technical paper draft.
%% The siggraph format is obsolete in current acmart; use format=acmtog,
%% which is the format used for SIGGRAPH technical papers and TOG.
%% Compile with:  pdflatex main -> bibtex main -> pdflatex main -> pdflatex main
\documentclass[acmtog, review=true, authorversion=false]{acmart}

%% --- Packages ---
\usepackage{graphicx}
\usepackage{amsmath}
% Note: do not load amssymb under acmart — newtxmath already defines \Bbbk.
\usepackage{booktabs}
\usepackage{xcolor}
\usepackage{subcaption}
\usepackage{listings}
\usepackage{hyperref}

%% --- Custom math macros ---
\newcommand{\vect}[1]{\mathbf{#1}}
\newcommand{\mat}[1]{\mathbf{#1}}
\DeclareMathOperator{\ASG}{ASG}
\DeclareMathOperator{\GGX}{GGX}
\DeclareMathOperator{\BRDF}{BRDF}
\DeclareMathOperator{\NDF}{NDF}

%% --- ACM metadata: replace with real values before submission ---
\acmJournal{TOG}
\acmVolume{0}
\acmNumber{0}
\acmArticle{0}
\acmMonth{0}
\acmYear{2026}

\begin{document}

\title{Anisotropic Gaussian Splatting with Normal-Mapped Microfacet Shading}

\author{Anonymous Author}
\email{anonymous@example.com}
\affiliation{%
  \institution{Anonymous Institution}
  \city{City}
  \country{Country}
}

\begin{abstract}
We present a method that combines normal-map textures with anisotropic
Gaussian splatting to enable view-consistent, relightable rendering of
surfaces exhibiting anisotropic reflectance (brushed metal, hair, fabric).
Our approach augments each 2D Gaussian primitive with a tangent-space
normal perturbation and shades it with an Anisotropic Spherical Gaussian
(ASG) microfacet distribution under a deferred rendering pipeline. Unlike
prior work that treats anisotropic appearance as a learned view-dependent
field, we derive a physically grounded BRDF in which the ASG long axis is
aligned to the Gaussian primitive's tangent, yielding sharper and more
stable anisotropic highlights while remaining compatible with real-time
splatting.
\end{abstract}

\ccsdesc[500]{Computing methodologies~Rendering}
\ccsdesc[300]{Computing methodologies~Reflectance modeling}

\keywords{Gaussian splatting, anisotropic BRDF, normal mapping, inverse
rendering, relighting}

\maketitle

% =============================================================
\section{Introduction}\label{sec:intro}

3D Gaussian Splatting (3DGS)~\cite{kerbl3Dgaussians} has become a
dominant primitive for real-time novel-view synthesis thanks to its
explicit, highly parallel rasterization. However, vanilla 3DGS models
view-dependent appearance with Spherical Harmonics (SH) and lacks an
explicit surface normal, which prevents physically-based shading,
relighting, and---critically---anisotropic reflectance effects such as
brushed metal or hair.

Two threads of work address these limitations. The first replaces SH
with an \emph{Anisotropic Spherical Gaussian} (ASG) appearance
field~\cite{specgaussian,glossgau}, directly modeling anisotropic
view-dependent color but without a physical BRDF decomposition. The
second augments Gaussians with normals and BRDF parameters for
inverse rendering~\cite{gsir,relightable3dgaussians,deferredgs}, enabling
relighting but relying on isotropic microfacet models.

In this paper we bridge the two: we attach a tangent-space normal
texture to each Gaussian primitive and shade with an ASG microfacet
NDF whose long axis is aligned to the primitive tangent. This yields
the first (to our knowledge) Gaussian splatting representation that
jointly supports normal-map detail and physical anisotropic BRDF
shading.

\paragraph{Contributions.}
\begin{itemize}
  \item A per-Gaussian tangent-space normal perturbation field that
        adds fine surface detail without changing primitive geometry.
  \item An anisotropic microfacet shading model for Gaussian splatting
        based on ASG NDF aligned to the primitive tangent.
  \item A deferred splatting pipeline that combines the two for
        real-time, relightable, anisotropic rendering.
\end{itemize}

% =============================================================
\section{Related Work}\label{sec:related}

\subsection{Gaussian Splatting}
Kerbl et al.~\cite{kerbl3Dgaussians} introduced 3D Gaussian Splatting
as an explicit, real-time radiance field representation. Subsequent
work improved geometry with 2D Gaussian Splatting~\cite{2dgaussians}
and surface-grounded primitives~\cite{geosplatting}.

\subsection{Anisotropic Appearance Modeling}
Spec-Gaussian~\cite{specgaussian} replaces SH with an ASG appearance
field for view-dependent, anisotropic color. GlossGau~\cite{glossgau}
uses an ASG to approximate the microfacet NDF for glossy inverse
rendering. SpecGaussian-LF~\cite{specgaussianlf} further improves
view-dependent modeling with latent features. GOGS~\cite{gogs}
captures anisotropic highlights on Gaussian surfels, while
SSD-GS~\cite{ssdgs} uses an anisotropic Fresnel-based model for
scattering and shadow decomposition. The ASG itself originates from
the work of Xu et al.~\cite{xu2013anisg}, and the Cook--Torrance
microfacet framework~\cite{cooktorrance,walter2007microfacet} and
the Disney principled BRDF~\cite{burley2012disney} provide the
broader physically-based shading context.

\subsection{Relightable Gaussian Splatting}
GS-IR~\cite{gsir} introduced depth-derivation normal regularization
for inverse rendering. Relightable 3D
Gaussians~\cite{relightable3dgaussians} attaches normals and BRDF
parameters per primitive with point-based ray tracing.
DeferredGS~\cite{deferredgs} decouples texture and lighting via
deferred shading with SDF-distilled normals. RGS-DR~\cite{rgsdr}
adds specular residuals in a 2DGS deferred pipeline. Spec-Gloss
Surfels~\cite{specglosssurfels} integrates a microfacet BRDF into
2DGS with diffusion-based normal priors. GI-GS~\cite{gigs} performs
global-illumination decomposition with deferred shading and path
tracing, GUS-IR~\cite{gusir} unifies forward and deferred shading,
RTR-GS~\cite{rtrgs} combines radiance transfer with deferred
reflections, and 3DSS~\cite{surf3dss} introduces the first
differentiable surface-splatting renderer for PBR inverse rendering.
The tangent-space normal-perturbation pattern we adopt follows
NVDIFFREC~\cite{nvdiffrec}, and the SG lighting and prefiltered
environment map follow~\cite{wang2009sphericalgaussian}.

\subsection{Per-Primitive Spatially-Varying Materials}
TextureSplat~\cite{texturesplat} attaches per-primitive texture maps
for spatially-varying normals and materials, and
SVG-IR~\cite{svgir} proposes per-Gaussian spatially varying
parameters. Our normal perturbation field is complementary: it is
specifically aligned to the ASG tangent frame so that the
perturbation direction modulates the anisotropic highlight pattern,
which isotropic NDFs cannot exploit.

\subsection{Normal Mapping in Neural Rendering}
Classical normal mapping perturbs the surface normal in tangent space
to add fine detail without tessellation. In neural rendering, normal
priors from monocular estimators~\cite{rosgs} and photometric
stereo~\cite{psgs} have been used to regularize Gaussian geometry.
We instead use a normal texture for \emph{shading} detail, leaving
geometry stable.

% =============================================================
\section{Method}\label{sec:method}

\subsection{Preliminaries}
Each 2D Gaussian primitive carries a center $\vect{c}\in\mathbb{R}^3$,
two tangent vectors $\vect{u},\vect{v}$ defining the disk, opacity
$\alpha$, and base color. Its surface normal is
$\vect{n}_0 = \vect{u}\times\vect{v}$.

\subsection{Per-Gaussian Normal Texture}
We predict a tangent-space normal perturbation
$\vect{d}\in[-1,1]^2$ per primitive from a small MLP conditioned on
the primitive's spatial features. The shaded normal is
\[
  \vect{n} = \mathrm{normalize}\!\left(
    \vect{n}_0 + d_x\,\vect{u} + d_y\,\vect{v}
  \right).
\]

\subsection{Anisotropic ASG Microfacet Shading}
We model the BRDF with an anisotropic NDF given by an ASG:
\[
  D(\vect{h}) = \ASG(\vect{h};\vect{n},\vect{t},\vect{b},\alpha_t,\alpha_b)
  = \exp\!\left(
      -\left(\frac{(\vect{h}\cdot\vect{t})^2}{\alpha_t^2}
            +\frac{(\vect{h}\cdot\vect{b})^2}{\alpha_b^2}\right)
    \right),
\]
where $\vect{h}$ is the half vector, $\vect{t}$ the tangent (aligned to
the Gaussian long axis), $\vect{b}=\vect{n}\times\vect{t}$ the bitangent,
and $\alpha_t,\alpha_b$ control anisotropy. The long axis of the ASG
is thus geometrically tied to the primitive shape.

\subsection{Deferred Splatting}
We splat a G-buffer (albedo, $\vect{n}$, $\vect{t}$, $\alpha_t$,
$\alpha_b$, depth) and shade in pixel space, following the deferred
design of DeferredGS~\cite{deferredgs} and
RGS-DR~\cite{rgsdr}.

% =============================================================
\section{Results}\label{sec:results}

Placeholder for experiments, comparisons against
Spec-Gaussian~\cite{specgaussian}, GlossGau~\cite{glossgau},
Relightable 3D Gaussians~\cite{relightable3dgaussians}, and
DeferredGS~\cite{deferredgs}.

% =============================================================
\section{Conclusion}\label{sec:conclusion}

We presented an anisotropic Gaussian splatting representation with
normal-mapped microfacet shading. By tying the ASG NDF to the
primitive tangent and perturbing the normal via a per-primitive
texture, our method achieves sharp, stable anisotropic highlights
while remaining relightable.

\bibliographystyle{ACM-Reference-Format}
\bibliography{references}

\end{document}
